\section{Experimental Setup}
\label{sec:setup}

\subsection{Datasets and preprocessing}
\label{subsec:setup-data}

We evaluate on five U.S.-state splits of the public \emph{Gowalla} check-in dataset \cite{cho2011gowalla}: AL, AZ, FL, CA, TX. Each split is filtered to users with at least five check-ins, applied per state before fold creation so each state preserves its native cardinality. Per-state statistics are in Table~\ref{tab:datasets}; FL/CA/TX serve as the headline scale and AL/AZ as smaller-scale anchors. Region labels come from the per-state region partition used by the substrate engine.

\begin{table}[t]
\centering
\caption{Dataset statistics for the five Gowalla \cite{cho2011gowalla} U.S.-state splits, after filtering to users with at least five check-ins. The headline three (FL/CA/TX) sit above the rule; the smaller-scale anchors (AL/AZ) sit below. Region cardinality is the prediction target on the next-region task.}
\label{tab:datasets}
\footnotesize
\renewcommand{\arraystretch}{1.1}
\begin{tabular}{l r r r r r}
\toprule
State & Users & Check-ins & POIs & Regions & Mean traj.\ length \\
\midrule
Florida (FL)    & 13{,}935 & 1{,}392{,}262 &  76{,}266 & 4{,}703 &  99.9 \\
California (CA) & 26{,}171 & 3{,}148{,}594 & 168{,}771 & 8{,}501 & 120.3 \\
Texas (TX)      & 27{,}603 & 4{,}067{,}543 & 160{,}575 & 6{,}553 & 147.4 \\
\midrule
Alabama (AL)    &  1{,}622 &   109{,}695   &  11{,}666 & 1{,}109 &  67.6 \\
Arizona (AZ)    &  3{,}331 &   227{,}956   &  20{,}469 & 1{,}547 &  68.4 \\
\bottomrule
\end{tabular}
\end{table}


For each user we form non-overlapping length-9 windows over the chronological sequence and predict the category and region of the tenth check-in (shorter sequences padded). Folds are \textbf{user-disjoint} via \texttt{StratifiedGroupKFold} with five folds (random\_state=42), so train/val/test never share a user; splits are reused across substrates, heads, and recipes so paired tests are well defined. The next-region transition prior is rebuilt per fold from training-fold edges only and seed-tagged. Check2HGI and HGI are pre-trained on each state's full check-in graph; the user-disjoint protocol applies to the downstream classifier and the per-fold transition prior but not to the substrate's exposure, so the setup is transductive at the substrate level.

\subsection{Baselines and metrics}
\label{subsec:setup-baselines}

The headline contrast is the cross-attention multi-task model against matched-head single-task (STL) ceilings: a GRU category head and a STAN-based region head with an additive log-transition prior over training-fold-only edges~\cite{luo2021stan}. \textbf{The same head classes appear on both sides of the contrast, so any difference between the multi-task and single-task rows is attributable to the sharing architecture, not to head capacity.} External baselines are POI-RGNN and MHA+PE (category) and Markov-1-region, single-task STAN, and ReHDM (region); we report them in \S\ref{sec:results}.

Category metrics: macro-F1 (primary) and Acc@1/Acc@3 (top-$k$ accuracy, secondary, since the 7-way distribution is skewed). Region metrics: Acc@10 (primary), Acc@5 and MRR (mean reciprocal rank, secondary). The joint score is $\Delta_m$~\cite{maninis2019attentive,vandenhende2022mtl} against the matched-head STL ceiling, reported on a primary pair (macro-F1, MRR) and a secondary pair (macro-F1, Acc@10).

\subsection{Statistical protocol}
\label{subsec:setup-stats}

Paired contrasts use the paired Wilcoxon signed-rank test with \texttt{alternative='greater'} for directional claims and the two-one-sided test (\textbf{TOST}) for non-inferiority at $\delta\in\{2,3\}$ pp. At the single-seed five-fold scale the max significance for $n=5$ is $p=0.0312$ (all folds agreeing).

Two ceiling protocols sit side by side. The \textbf{substrate axis} (Table~\ref{tab:substrate}) uses the seed-42 five-fold paired Wilcoxon at the $n=5$ ceiling; multi-seed mean and seed-$\sigma$ are shown for Check2HGI at AL/AZ/FL but the paired test stays single-seed. The \textbf{architectural axis} (Tables~\ref{tab:mtl-stl}, \ref{tab:external}) pools four extra seeds $\{0,1,7,100\}$ on top of seed-42, yielding $n=20$ fold-pairs per state; FL's $\Delta_m$ pair additionally carries $n=25$ fold-pairs. Effect sizes are in pp for F1, Acc@$k$, and MRR; $\Delta_m$ is in percent.
